\documentclass{beamer}
\usepackage[utf8]{inputenc}

\usetheme{Madrid}
\usecolortheme{default}
\usepackage{amsmath,amssymb,amsfonts,amsthm}
\usepackage{txfonts}
\usepackage{tkz-euclide}
\usepackage{listings}
\usepackage{adjustbox}
\usepackage{array}
\usepackage{tabularx}
\usepackage{gvv}
\usepackage{lmodern}
\usepackage{circuitikz}
\usepackage{tikz}
\usepackage{graphicx}

\setbeamertemplate{page number in head/foot}[totalframenumber]

\usepackage{tcolorbox}
\tcbuselibrary{minted,breakable,xparse,skins}



\definecolor{bg}{gray}{0.95}
\DeclareTCBListing{mintedbox}{O{}m!O{}}{%
	breakable=true,
	listing engine=minted,
	listing only,
	minted language=#2,
	minted style=default,
	minted options={%
		linenos,
		gobble=0,
		breaklines=true,
		breakafter=,,
		fontsize=\small,
		numbersep=8pt,
		#1},
	boxsep=0pt,
	left skip=0pt,
	right skip=0pt,
	left=25pt,
	right=0pt,
	top=3pt,
	bottom=3pt,
	arc=5pt,
	leftrule=0pt,
	rightrule=0pt,
	bottomrule=2pt,
	toprule=2pt,
	colback=bg,
	colframe=orange!70,
	enhanced,
	overlay={%
		\begin{tcbclipinterior}
			\fill[orange!20!white] (frame.south west) rectangle ([xshift=20pt]frame.north west);
	\end{tcbclipinterior}},
	#3,
}
\lstset{
	language=C,
	basicstyle=\ttfamily\small,
	keywordstyle=\color{blue},
	stringstyle=\color{orange},
	commentstyle=\color{green!60!black},
	numbers=left,
	numberstyle=\tiny\color{gray},
	breaklines=true,
	showstringspaces=false,
}
%------------------------------------------------------------
%This block of code defines the information to appear in the
%Title page
\title %optional
{5.2.49}
%\subtitle{A short story}

\author % (optional)
{RAVULA SHASHANK REDDY - EE25BTECH11047}

 \begin{document}
	
	
	\frame{\titlepage}
	\begin{frame}{Question}
    The eigenvalues of the matrix 
\[
\vec{A} = \myvec{3 & -1 & 1 \\ -1 & 5 & -1 \\ 1 & -1 & 3}
\]

are $\lambda_1, \lambda_2$ and $\lambda_3$. Find the value of 
\[
\lambda_1 \lambda_2 \lambda_3 (\lambda_1 + \lambda_2 + \lambda_3).
\]
\end{frame}
\begin{frame}{Solution}
    We know that for any square matrix:
\begin{align}
\lambda_1 \lambda_2 \lambda_3 &= \det \vec{A}, \qquad 
\lambda_1 + \lambda_2 + \lambda_3 = \texr{trace} (\vec{A}).\\
 \text{trace}(\vec{A}) &= 3 + 5 + 3 = 11.
\end{align}
\begin{align}
\det \vec{A} 
&= 3(15-1) - (-1)(-3+1) + (1-5) \\[1mm]
&= 3\cdot 14 - 2 - 4 \\[1mm]
&= 36.
\end{align}

Hence
\begin{align}
\lambda_1 \lambda_2 \lambda_3 (\lambda_1 + \lambda_2 + \lambda_3) = \det \vec{A} \cdot \text{trace} (\vec{A}) = 36 \cdot 11 = 396.
\end{align}
\end{frame}

\begin{frame}[fragile]
\frametitle{C Code}
\begin{lstlisting}
    #include <stdio.h>

int main() {
    int A[3][3] = {
        {3, -1, 1},
        {-1, 5, -1},
        {1, -1, 3}
    };

    int trace = 0;
    for (int i = 0; i < 3; i++) {
        trace += A[i][i];
    }

    // Compute determinant using cofactor expansion along first row
    \end{lstlisting}
    \end{frame}
\begin{frame}[fragile]
\frametitle{C Code}
\begin{lstlisting}
    int det = 0;
    for (int j = 0; j < 3; j++) {
        int minor = A[1][(j+1)%3]*A[2][(j+2)%3] - A[1][(j+2)%3]*A[2][(j+1)%3];
        if (j % 2 == 0)
            det += A[0][j] * minor;
        else
            det -= A[0][j] * minor;
    }

    int result = det * trace;

    printf("Trace = %d\n", trace);
    printf("Determinant = %d\n", det);
    printf("lambda1*lambda2*lambda3*(lambda1+lambda2+lambda3) = %d\n", result);

    return 0;
}
\end{lstlisting}
\end{frame}
\begin{frame}[fragile]
\frametitle{Python Direct}
\begin{lstlisting}
    import numpy as np

# Define the matrix
A = np.array([[3, -1, 1],
              [-1, 5, -1],
              [1, -1, 3]])

trace = np.trace(A)

# Compute determinant
det = np.linalg.det(A)

result = det * trace

print("Trace =", trace)
print("Determinant =", det)
print("lambda1*lambda2*lambda3*(lambda1+lambda2+lambda3) =", result)
\end{lstlisting}
\end{frame}
\begin{frame}[fragile]
\frametitle{Python Shared}
\begin{lstlisting}
    import ctypes

# Load the shared library
lib = ctypes.CDLL('./libmatrix_calc.so')

# Tell ctypes the return type of the function
lib.compute_result.restype = ctypes.c_int

# Call the function
result = lib.compute_result()

print("lambda1*lambda2*lambda3*(lambda1+lambda2+lambda3) =", result)
\end{lstlisting}
\end{frame}
\end{document}
